\documentclass[../main.tex]{subfiles}

\begin{document}

\chapter{LDPC Codes: Theory and Applications}
% Sources: 
%       Gallager, R. G. (1962). Low-Density Parity-Check Codes. IRE Transactions on Information Theory, 8(1), 21–28. DOI:10.1109/TIT.1962.1057683
An LDPC code is a linear block code defined by a sparse binary parity-check matrix $H \in \mathbb{F}_2^{m \times n}$, where:
\begin{itemize}
    \item $n$: length of the codeword
    \item $m$: number of parity-check equations
    \item $\mathbb{F}_2$: the finite field with two elements {0, 1}.
\end{itemize}

A vector $c \in \mathbb{F}_2^{n}$ is a valid codeword if and only if:
$$ H \cdot c^T = 0 mod 2 $$

This condition implies that every codeword lies in the null space (kernel) of $H$, over GF(2). Since $H$ has rank  $r \leq m$, the dimension of the code is $k = n - rank(H)$, and the code rate is:
$$R = \frac{k}{n} = 1 - \frac{rank(H)}{n}$$

% ########## ENCODING ##########
% Sources:
%       Richardson, T., & Urbanke, R. (2008). Modern Coding Theory. Cambridge University Press. (Chapter 4)
\section{LDPC Encoding}
\label{sec:encoding}
While decoding is typically based on the parity-check matrix $H$, encoding is performed using a generator matrix $G \in \mathbb{F}_2^{k \times n}$, which satisfies:
$$G \cdot H^T = 0$$
This ensures that any message vector $u \in \mathbb{F}_2^{k}$ is mapped to a valid codeword $c \in \mathbb{F}_2^{n}$:
$$c = u \cdot G$$
To construct $G$ from $H$, one common approach is to transform $H$ into systematic form using Gaussian elimination:
$$
H = \begin{bmatrix} P & I \end{bmatrix}
\quad \Rightarrow \quad
G = \begin{bmatrix} I & P^T \end{bmatrix}
$$
This allows the parity bits to be computed from the message bits efficiently.

% ########## DECODING ##########
% Sources:
\section{LDPC Decoding}
\label{sec:decoding}
LDPC decoding is typically performed using iterative algorithms on dhe code's Tanner graph, These algorithms aim to estimate the transmitted codeword by exploiting the sparsity of the parity-check matrix $H$ and the soft information received from the channel. 
This section presents the mathematical formulation of the Sum-Product Algorithm (SPA), as well as its approximations: Min-Sum, Normalized Min-Sum, and Offset Min-Sum.

% ########## DECODING ALGORITHMS ##########
\section{Overview Decoding Algorithms}
\label{sec:decoding_algorithms}

\subfile{./4-theory/SPA.tex}

\subfile {./4-theory/MSA.tex}

\subfile{./4-theory/NMS-OMS.tex}

\end{document}